\chapter{BRCA Genetic Testing}

\section*{What Is This Test?}
BRCA1 and BRCA2 genetic testing analyzes the BRCA1 gene (on chromosome 17q21) and BRCA2 gene (on chromosome 13q13.1) for germline pathogenic or likely pathogenic variants that confer a markedly increased lifetime risk of breast, ovarian, prostate, pancreatic, and other cancers. The BRCA1 and BRCA2 proteins are critical tumor suppressors that function in the homologous recombination repair (HRR) pathway to repair DNA double-strand breaks. Loss-of-function germline mutations in BRCA1 or BRCA2 result in defective homologous recombination, genomic instability, and a predisposition to cancer. Women with a pathogenic BRCA1 variant have a 55--72\% lifetime risk of breast cancer and a 39--44\% lifetime risk of ovarian cancer; BRCA2 carriers have a 45--69\% lifetime breast cancer risk and an 11--17\% ovarian cancer risk \cite{kuchenbaecker2017risks}. Men with BRCA2 have an elevated risk of male breast cancer (7\% lifetime) and prostate cancer (20--25\%). Both genes also confer increased risks for pancreatic cancer and melanoma. Testing is performed by next-generation sequencing (NGS) with Sanger confirmation of pathogenic variants, and deletion/duplication analysis (MLPA or chromosomal microarray) to detect large gene rearrangements. Identification of a pathogenic variant guides cancer risk management (enhanced screening, risk-reducing surgery, chemoprevention) and identifies family members who may carry the same variant (cascade testing).

\section*{Alternative Names}
BRCA1/BRCA2 Testing; Hereditary Breast and Ovarian Cancer Syndrome Testing; BRCA Mutation Analysis; BRCA Gene Sequencing; HBOC Testing; Breast Cancer Gene Testing

\section*{Principle}
Germline BRCA1 and BRCA2 testing is performed on genomic DNA extracted from peripheral blood leukocytes (or saliva). Full gene sequencing of the coding exons and flanking intronic splice sites of BRCA1 (24 exons) and BRCA2 (27 exons) is performed by next-generation sequencing (NGS) using targeted capture and massively parallel sequencing. Alignment to the reference genome (GRCh38/hg38) and variant calling are performed using bioinformatics pipelines. Variants are classified according to the American College of Medical Genetics and Genomics (ACMG) and the Evidence-Based Network for the Interpretation of Germline Mutant Alleles (ENIGMA) consortium criteria into five categories: pathogenic (Class 5), likely pathogenic (Class 4), variant of uncertain significance (VUS, Class 3), likely benign (Class 2), and benign (Class 1) \cite{richards2015standards}. Large gene rearrangements (deletions and duplications of one or more exons), which account for 5--15\% of all BRCA mutations, are detected by multiplex ligation-dependent probe amplification (MLPA) or chromosomal microarray analysis (CMA). Results are reported as: positive (pathogenic or likely pathogenic variant identified), negative (no pathogenic variant identified), or VUS (variant of uncertain significance identified).

\section*{History}
The BRCA1 gene was localized to chromosome 17q21 by Mary-Claire King's laboratory in 1990 \cite{hall1990linkage} and cloned by a collaborative effort led by the University of Utah, Myriad Genetics, and NIH in 1994 \cite{miki1994strong}. BRCA2 was localized to chromosome 13q12--13 in 1994 \cite{wooster1994localization} and cloned in 1995 by the Institute of Cancer Research (UK) and collaborating groups \cite{wooster1995identification}. Myriad Genetics obtained patents on the BRCA1 and BRCA2 genes in the 1990s, making it the exclusive provider of clinical BRCA testing in the United States from 1996 to 2013. In the landmark 2013 Supreme Court case \textit{Association for Molecular Pathology v. Myriad Genetics, Inc.}, the Court unanimously ruled that naturally occurring DNA sequences cannot be patented, opening BRCA testing to multiple laboratories and dramatically reducing costs \cite{myriad2013supremecourt}. The advent of NGS has enabled the expansion from single-gene BRCA testing to multigene hereditary cancer panels (testing 20--80+ cancer susceptibility genes simultaneously). In 2014 and 2018, the FDA approved the PARP inhibitors olaparib, rucaparib, niraparib, and talazoparib for BRCA-associated ovarian, breast, pancreatic, and prostate cancers, leveraging synthetic lethality: cancer cells with BRCA mutations rely on PARP-mediated DNA repair, and PARP inhibition causes catastrophic DNA damage and cell death in BRCA-deficient tumors \cite{robson2017olaparib}.

\section*{Why This Test Is Ordered}
NCCN guidelines recommend BRCA testing for: (1) personal history of breast cancer diagnosed at age $\leq$45, (2) personal history of triple-negative breast cancer at age $\leq$60, (3) male breast cancer at any age, (4) ovarian cancer (including fallopian tube and primary peritoneal) at any age, (5) pancreatic cancer at any age, (6) metastatic prostate cancer at any age, (7) two or more primary breast cancers in the same individual, (8) Ashkenazi Jewish ancestry with breast or pancreatic cancer, (9) family history of known BRCA1/BRCA2 pathogenic variant, and (10) family history suggestive of hereditary breast and ovarian cancer syndrome (multiple affected relatives) \cite{nccn2024genetic}. The purpose is to identify individuals with hereditary cancer predisposition who will benefit from risk-reducing interventions and to guide PARP inhibitor therapy selection in patients with BRCA1/BRCA2-associated cancers.

\section*{Primary vs Secondary Test}
This is a primary genetic test for hereditary cancer predisposition. For individuals who meet NCCN criteria, BRCA testing is the standard of care for cancer risk assessment and management. In the context of ovarian, pancreatic, and metastatic prostate cancer, BRCA testing also serves as a companion diagnostic for PARP inhibitor eligibility.

\section*{How This Test Is Performed}
For most patients, a health care professional will take a blood sample from a vein in your arm using a small needle. After the needle is inserted, a small amount of blood will be collected into a test tube or vial. You may feel a little sting when the needle goes in or out. This usually takes less than five minutes. Alternatively, saliva collected into an Oragene DNA self-collection kit may be used. Testing is ideally performed after pre-test genetic counseling to discuss the implications, limitations, and potential outcomes of the test.

\section*{What Preparation Is Needed}
No special preparation is required. Fasting is not necessary. Patients should undergo pre-test genetic counseling with a certified genetic counselor or a healthcare provider experienced in cancer genetics. Testing is ideally performed when the patient is not acutely ill. Allogeneic hematopoietic stem cell transplant recipients should not have blood-based genetic testing, as the DNA in circulating leukocytes is donor-derived; a pre-transplant blood sample or an alternative tissue (skin biopsy for fibroblast culture, buccal swab) should be used.

\section*{Contraindications and Precautions}
The primary contraindication is the absence of informed consent and pre-test genetic counseling. Genetic testing has profound implications for the patient and their family members. The Genetic Information Nondiscrimination Act (GINA) of 2008 protects against genetic discrimination in health insurance and employment but does NOT apply to life insurance, disability insurance, or long-term care insurance. Patients should be counseled about potential implications for insurance before testing. A variant of uncertain significance (VUS) is a common result (reported in 5--10\% of tests); a VUS is NOT actionable and should NOT be used for clinical decision-making. VUS classification may change over time as more data accumulate; the laboratory may reclassify the variant and issue an updated report. Parents should be advised that AR (autosomal recessive) Fanconi anemia type D1 (biallelic BRCA2) and Fanconi anemia type S (biallelic BRCA1) can occur in offspring if both parents are carriers. Specimen should be collected in an EDTA (lavender-top) tube for DNA extraction. Heparin (green-top) tubes interfere with PCR-based assays.

\section*{Risks}
Minimal physical risk (bruising, pain at the venipuncture site). The greater risks are psychological (anxiety, depression, guilt) related to a positive result, and potential genetic discrimination in life, disability, or long-term care insurance (not covered by GINA). Family members of a person with a positive BRCA result may experience distress about their own risk.

\section*{What Happens After the Test}
Results are typically available in 2--4 weeks. Post-test genetic counseling should be provided to discuss the results and their implications. A positive result (pathogenic or likely pathogenic variant identified) triggers specific risk management recommendations: breast MRI with contrast annually starting at age 25--30 and annual mammography starting at age 30, discussion of risk-reducing bilateral salpingo-oophorectomy (RRSO) at age 35--40 for BRCA1 and 40--45 for BRCA2, discussion of risk-reducing mastectomy, chemoprevention with tamoxifen or raloxifene, and cascade testing of at-risk family members (parents, siblings, children, who each have a 50\% probability of carrying the same variant). A negative result in the context of a known familial variant is a ``true negative'' and indicates that the individual has not inherited the familial mutation and their cancer risks are comparable to those of the general population. A negative result in the absence of a known familial variant must be interpreted with caution: if the family mutation was not identified (or if testing was not comprehensive), a negative result does not rule out an inherited cancer predisposition.

\section*{Understanding Results}

\subsection*{Pathogenic/Likely Pathogenic Variant Detected (Positive)}
NCCN cancer risk management guidelines include: (1) Breast awareness starting at age 18, clinical breast examination every 6--12 months starting at age 25. (2) Annual breast MRI with contrast starting at age 25--29 (BRCA1) or 25--29 (BRCA2), and annual mammogram and breast MRI starting at age 30--75. (3) Discussion of risk-reducing bilateral mastectomy. (4) Risk-reducing bilateral salpingo-oophorectomy (RRSO) at age 35--40 (BRCA1) or 40--45 (BRCA2) upon completion of childbearing. The lifetime ovarian cancer risk reduction with RRSO is 80--90\% \cite{domchek2010association}. (5) For BRCA2 carriers: prostate cancer screening starting at age 40 (PSA and DRE). (6) Pancreatic cancer screening (annual MRI/MRCP or EUS) starting at age 50 or 10 years before the earliest pancreatic cancer in the family, if a first-degree relative with pancreatic cancer exists. (7) Cascade genetic testing of at-risk family members. (8) For patients with BRCA-associated cancers, PARP inhibitor therapy (olaparib for ovarian, breast, pancreatic, and prostate cancer; niraparib and rucaparib for ovarian cancer; talazoparib for breast cancer).

\subsection*{No Pathogenic Variant Detected (Negative)}
A negative BRCA test result must be interpreted in the clinical and family history context. If a known familial BRCA pathogenic variant has been identified and the patient tests negative for that variant, this is a true negative: the patient has NOT inherited the familial predisposition, and their cancer risks approximate those of the general population. Enhanced screening is not required. If no BRCA pathogenic variant is detected and no familial variant is known, the negative result does NOT exclude hereditary cancer predisposition. Other genes (PALB2, CHEK2, ATM, TP53, PTEN, STK11, CDH1) may harbor pathogenic variants. Family history-based risk-appropriate screening should continue. A negative BRCA result in an individual who has already had cancer does not rule out a somatic BRCA mutation in the tumor; somatic BRCA testing on tumor tissue may be warranted for PARP inhibitor eligibility.

\subsection*{Variant of Uncertain Significance (VUS)}
A VUS is a genetic change for which available evidence is insufficient to classify it as pathogenic or benign. Approximately 5--10\% of BRCA test results include a VUS. A VUS is NOT actionable and should NEVER be used for clinical decision-making (risk-reducing surgery should not be recommended based on a VUS alone). VUS reclassification over time is expected: most VUSs are eventually downgraded to likely benign/benign as more data are accumulated. Laboratories periodically issue updated reports when reclassification occurs. Patients should maintain contact with their genetic counselor for updates on VUS classification.

\section*{Normal Results}
\begin{itemize}
  \item \textbf{No pathogenic variant detected:} The most common result. Normal population frequency of pathogenic BRCA1/BRCA2 variants is approximately 1:300 to 1:500 (0.2--0.3\%) in the general non-Ashkenazi Jewish population. In the Ashkenazi Jewish population, approximately 1:40 (2.5\%) carry one of three founder mutations (BRCA1 185delAG, BRCA1 5382insC, BRCA2 6174delT)
  \item \textbf{Benign variant:} A genetic variant that has been confirmed NOT to increase cancer risk. Reported in supplemental data or not reported at all. No clinical action required
  \item \textbf{Likely benign variant:} Insufficient evidence to classify as definitively benign but unlikely to be pathogenic. No clinical action required
  \item \textbf{Variant of uncertain significance (VUS):} Neither pathogenic nor benign. No clinical action recommended. Follow-up for future variant reclassification
  \item Results must be interpreted with reference to the ACMG/AMP variant classification guidelines and the ENIGMA consortium classification rules
\end{itemize}

\section*{Follow-up Tests}
Cascade genetic testing of first-degree relatives for the identified familial pathogenic variant, single-site targeted testing (significantly less expensive than full gene sequencing). Multigene hereditary cancer panel testing (PALB2, CHEK2, ATM, RAD51C, RAD51D, BRIP1, TP53, PTEN, STK11, CDH1, NBN, NF1, MSH2, MLH1, MSH6, PMS2, EPCAM) if negative for BRCA but with a strong family history. Breast MRI with and without contrast, annual. Mammogram with tomosynthesis, annual, starting at age 30--35. Transvaginal ultrasound and CA-125, every 6--12 months, starting at age 30--35 in patients who decline or defer RRSO (limited sensitivity and specificity; RRSO remains the standard of care). MRI/MRCP or EUS for pancreatic cancer screening starting at age 50 (or 10 years before earliest pancreatic cancer in family). PSA and DRE for prostate cancer screening starting at age 40 (BRCA2 carriers). Tumor genomic testing: somatic BRCA1/BRCA2 testing on tumor tissue if PARP inhibitor therapy is being considered for patients who have a cancer diagnosis but a negative germline BRCA test.

\section*{References}
\bibliographystyle{unsrtnat}
\bibliography{references}
\clearpage
\section*{Regulatory Status and Authorization Date}
Regulatory status applies to the specific assay, instrument, manufacturer, and intended use rather than to the test name alone. No single approval date can be assigned to this test category. For a commercial assay, verify the current product in the relevant regulator's database: the U.S. Food and Drug Administration (FDA) clearance, approval, or authorization; India's Central Drugs Standard Control Organisation (CDSCO) licence or permission; the European Union In Vitro Diagnostic Regulation (EU IVDR) CE certificate issued through the applicable conformity-assessment route; or the United Kingdom Medicines and Healthcare products Regulatory Agency (MHRA) registration and UKCA/CE status. Laboratory-developed tests may not have an individual FDA, CDSCO, EU IVDR, or MHRA approval date.
\clearpage
